
\subsubsection{5.1 Main Results}

\textbf{Table 1: Clustering results for annotation-free spurious direction discovery at epoch 5, k=2.}

\begin{table}[htbp]
\centering
\begin{tabular}{llllllll}
\toprule
Dataset & AMI & Purity & AMI (random) & Purity (random) & AMI ≥ 0.5 & Purity ≥ 0.75 & Pass \\
\midrule
Waterbirds & 0.762 & 0.892 & 0.0001 & 0.728 & Yes & Yes & Pass \\
CelebA & 0.258 & 0.456 & 0.000 & 0.440 & No & No & Fail \\
\bottomrule
\end{tabular}
\end{table}

Waterbirds exceeds both thresholds: AMI is 0.262 above threshold and purity is 0.142 above threshold. CelebA fails both thresholds: AMI of 0.258 is 0.242 below the 0.5 threshold, and purity of 0.456 is 0.294 below the 0.75 threshold. The CelebA purity of 0.456 is 0.016 above the random baseline of 0.440, representing a margin that is close to chance. The overall gate result (MUST_WORK across both datasets) is FAIL.

Figure 1 presents a bar chart of AMI and worst-cluster purity for both datasets with threshold lines and random baselines indicated.

\textit{[Figure 1: figures/metrics_bar.png — AMI and worst-cluster purity for Waterbirds and CelebA with threshold lines (AMI=0.5, purity=0.75) and random baselines.]}

\subsubsection{5.2 t-SNE Visualizations}

t-SNE projections of epoch-5 embeddings are computed for both datasets and visualized with three colorings: k-means cluster assignment, ground-truth group identity, and class label.

For Waterbirds, the t-SNE projection colored by k-means cluster assignment shows two spatially separated regions corresponding to the two clusters. The t-SNE projection colored by ground-truth group identity shows a similar two-region structure with close correspondence to the cluster boundaries, consistent with the AMI of 0.762.

\textit{[Figure 2: figures/tsne_cluster_waterbirds.png — t-SNE of Waterbirds epoch-5 embeddings, colored by k-means cluster assignment.]}

\textit{[Figure 3: figures/tsne_group_waterbirds.png — t-SNE of Waterbirds epoch-5 embeddings, colored by ground-truth group identity.]}

For CelebA, the t-SNE projection colored by k-means cluster assignment does not exhibit a structure that aligns with the ground-truth group identities. Both spurious groups distribute across both clusters without spatially coherent separation, consistent with the near-random purity of 0.456.

\textit{[Figure 4: figures/tsne_cluster_celeba.png — t-SNE of CelebA epoch-5 embeddings, colored by k-means cluster assignment.]}

\subsubsection{5.3 Cluster Composition}

For Waterbirds (Figure 5), cluster 0 contains approximately 89\% water-background samples and cluster 1 contains approximately 89\% land-background samples. Each cluster is dominated by a single spurious group, with worst-cluster purity 0.892.

\textit{[Figure 5: figures/cluster_composition_waterbirds.png — Group composition within each k-means cluster for Waterbirds. Each cluster is dominated by one spurious group (purity=0.892).]}

For CelebA (Figure 6), both clusters contain mixed membership across all four spurious groups (non-blonde female, blonde female, non-blonde male, blonde male). No cluster is dominated by any single group. This pattern is consistent with the worst-cluster purity of 0.456 being near-random.

\textit{[Figure 6: figures/cluster_composition_celeba.png — Group composition within each k-means cluster for CelebA. Both clusters exhibit mixed membership across spurious groups (purity=0.456).]}

\subsubsection{5.4 Training State at Epoch 5}

\textbf{Table 2: Training loss and accuracy at epoch 5.}

\begin{table}[htbp]
\centering
\begin{tabular}{llll}
\toprule
Dataset & Train loss & Train accuracy & Embedding shape \\
\midrule
Waterbirds & 0.088 & 96.66\% & (4,795 × 2,048) \\
CelebA & 0.124 & 95.00\% & (162,770 × 2,048) \\
\bottomrule
\end{tabular}
\end{table}

Both models achieve high training accuracy at epoch 5. The CelebA model's training loss of 0.124 and accuracy of 95.00\% indicates that the model is learning the task and is not undertrained. The clustering failure on CelebA is not attributable to insufficient training progress at epoch 5.

\subsubsection{5.5 Comparison with Published Configurations}

Published annotation-free clustering results, including those attributed to PruSC \citep{Kimetal2024}, are reported under configurations that may differ substantially from the epoch-5, k=2 setting. Specifically, published results appear to be obtained at or near training convergence with k values chosen to match or exceed the number of known groups. The epoch-5, k=2 CelebA result of purity=0.456 (barely above the random baseline of 0.440) does not directly contradict convergence-epoch, k≥4 results; rather, it demonstrates that these two configurations are not interchangeable. Convergence-epoch results cannot be cited as evidence that early-epoch, k=2 annotation-free detection is reliable on CelebA. The configurations differ in ways that are practically significant for intervention methods that require detection at an early probe epoch, and these configuration differences are not disclosed in existing published reporting.

Note that the PruSC citation has not been independently verified through a literature search; see Section 6.3 Limitation L5 for discussion.

